\documentclass[12pt, reqno]{amsart}
\usepackage{amsmath, amsthm, amscd, amsfonts, amssymb, mathtools}
\usepackage{color, geometry, relsize, hyperref}
\usepackage[dvipsnames]{xcolor}
%\usepackage{amsaddr}
\textheight 22.5truecm \textwidth 14.5truecm
\setlength{\oddsidemargin}{0.35in}\setlength{\evensidemargin}{0.35in}

\setlength{\topmargin}{-.5cm}

\newtheorem{theorem}{Theorem}[section]
\newtheorem{lemma}[theorem]{Lemma}
\newtheorem{proposition}[theorem]{Proposition}
\newtheorem{corollary}[theorem]{Corollary}
\newtheorem{observation}[theorem]{Observation}
\newtheorem{fact}[theorem]{Fact}
\theoremstyle{definition}
\newtheorem{definition}[theorem]{Definition}
\newtheorem{example}[theorem]{Example}
\newtheorem{conjecture}[theorem]{Conjecture}
\newtheorem{solution}[theorem]{Solution}
\newtheorem{question}[theorem]{Question}
\theoremstyle{remark}
\newtheorem{remark}[theorem]{Remark}
\numberwithin{equation}{section}

\renewcommand{\baselinestretch}{1.5}



\begin{document}
\setcounter{page}{1}

\title[Eigenvalue bounds]
{Eigenvalue bounds for general matrix rational functions}
	
\author[Pallavi, Shrinath and Sachindranath]{Pallavi. B, Shrinath Hadimani 
and Sachindranath Jayaraman}
\address{School of Mathematics\\ 
Indian Institute of Science Education and Research Thiruvananthapuram\\ 
Maruthamala P.O., Vithura, Thiruvananthapuram -- 695 551, Kerala, India.}
\email{(pallavipoorna20, srinathsh3320, sachindranathj)@iisertvm.ac.in, 
sachindranathj@gmail.com}
	
	
\subjclass[2010]{15A18, 15A42, 47A12, 47A56, 26C15}
	
\keywords{Rational eigenvalue problems; matrix polynomials; matrix rational functions; 
bounds on eigenvalues; numerical radius}
	
\begin{abstract} 
Most nonlinear eigenvalue problems are rational eigenvalue 
problems or can be approximated to rational eigenvalue problems. Bounds 
on the eigenvalues of matrix rational functions (also known as rational matrices) 
are less studied in literature compared to matrix polynomials. Our aim is to derive 
bounds on the eigenvalues of an arbitrary matrix rational function in terms of 
its poles. A bound is obtained as a root of a real rational function associated 
with a matrix rational function. In addition, a block matrix associated with this 
rational function is used to get explicit bounds on its zeros and in turn on the 
eigenvalues of matrix rational function. Bounds on the zeros of an arbitrary complex 
rational function are also obtained in this process.
\end{abstract}
	
\maketitle
	
\section{Introduction}\label{sec-1}

The nonlinear eigenvalue problem (abbreviated as NEP) seeks to find a scalar $\lambda_0$ 
and a nonzero vector $v \in \mathbb{C}^n$ satisfying $G(\lambda_0)v=0$, where $G(\lambda)$ 
is a regular matrix-valued function (see Section \ref{sec-2} for definition), and each 
entry of $G(\lambda_0)$ is bounded. In many nonlinear eigenvalue problems which arise in 
applications, $G(\lambda)$ is a meromorphic function and hence is either a 
matrix rational function or can be approximated by a matrix rational function using
rational approximation methods \cite{Saad-Guide-Meidlar}. By a matrix rational function, 
we mean matrices whose entries are scalar rational functions. These are sometimes 
referred to as rational matrices in the literature. The focus on such matrices began in 
connection with study of their structural data, such as its poles, zeros and their partial 
multiplicities in linear systems and control theory \cite{Kailath}. 
A thorough study on NEPs can be found in \cite{Dopico}, \cite{Mehrmann-Voss} and 
the references cited therein. Given a matrix rational function $G(\lambda)$, one can 
determine the eigenvalues using iterative methods \cite{Mehrmann-Voss}, where one starts 
with an initial guess for the eigenvalue. Naturally, a {\it good guess} will reduce the 
number of iterative steps. There is, however, a different method where one goes through a  
rational approximation, where the eigenvalues are computed inside a closed 
contour. In both cases, determining the location of the eigenvalues is necessary 
to reduce the computational effort. 

In \cite{Pallavi}, the authors consider a particular type of matrix rational function 
and derive bounds on the eigenvalues using the Bauer-Fike theorem on an associated block 
matrix. Bounds were also derived by associating a real rational function to the given 
matrix rational function. The current work is an extension of \cite{Pallavi} 
for a more general form of matrix rational function. In particular, we derive bounds on 
the eigenvalues of an arbitrary matrix rational function by associating a block matrix 
as well as by associating a scalar rational function to the given matrix rational function. 
We use bounds on the numerical radius of a block matrix associated with the scalar rational function to derive bounds on the eigenvalues of a matrix rational function. The manuscript 
is organized as follows. Section \ref{sec-1} is introductory. Section \ref{sec-2} collects 
basic facts about matrix rational functions and inequalities involving the numerical 
radius that will be used later on. The main results are presented in Section \ref{sec-3}, 
which is subdivided into subsections for ease of reading. These subsections are 
self-explanatory. The manuscript ends with numerical examples illustrating the results 
obtained. 

\medskip
\section{Preliminaries}\hspace*{\fill}
\label{sec-2}

The space of all $p \times s$ matrices over complex numbers is denoted by $M_{p,s}$. 
When $p=s$, this will be denoted by $M_p$. For $A \in M_p$, the notations $||A||_2, 
||A||_\infty$ and $||A||_1$ denote respectively the spectral norm, the maximum row sum norm 
and the maximum column sum norm of $A$. A matrix-valued function $P(\lambda)= C_m \lambda^m + C_{m-1}\lambda^{m-1} + \cdots +C_0$, where $C_i \in M_p$ for $i=1, 2, \ldots, m$ and 
$C_m \neq 0$ is called a $p \times p$ matrix polynomial of degree $m$. A matrix rational 
function of size $p \times s$ is a matrix-valued function of the form $R(\lambda)=[r_{ij}(\lambda)]$, ($1 \leq i \leq p, 1 \leq j \leq s$) where 
$\displaystyle r_{ij}(\lambda) = \frac{p_{ij}(\lambda)}{q_{ij}(\lambda)}$ is a scalar 
rational function and $p_{ij}, q_{ij}$ are scalar polynomials. A square matrix rational 
function $R(\lambda)$ is said to be regular if the determinant of $R(\lambda)$ does not 
vanish identically. For a regular matrix rational function $R(\lambda)$ a scalar 
$\lambda_0$ is called an eigenvalue of $R(\lambda)$ if there exists a nonzero vector 
$v \in \mathbb{C}^n$ such that $R(\lambda_0)v=0$, with $R(\lambda_0)$ bounded. The vector 
$v$ so obtained is called an eigenvector of $R(\lambda)$ corresponding to the eigenvalue 
$\lambda_0$. 

Given a $p \times p$ matrix rational function $R(\lambda)$, let $d(\lambda)$ be the monic 
least common denominator of the entries of $R(\lambda)$. Then there exist unimodular 
matrix polynomials (square matrix polynomials with nonzero constant determinant) 
$U(\lambda)$ and $V(\lambda)$ of size $p\times p$ such that 
\begin{equation} \label{Eqn-Smith-McMillan form}
U(\lambda) R(\lambda) V(\lambda) = \begin{bmatrix}
\displaystyle \text{diag}\bigg(\frac{\epsilon_1(\lambda)}{\psi_1(\lambda)},\cdots,
\frac{\epsilon_r(\lambda)}{\psi_r(\lambda)}\bigg) & 0 \\
0 & 0
\end{bmatrix},
\end{equation}
where $r \leq p$, $\epsilon_i(\lambda), \psi_i(\lambda)$ are monic scalar polynomials 
with 
$\displaystyle \frac{\epsilon_i(\lambda)}{\psi_i(\lambda)}$ irreducible and 
$\epsilon_i(\lambda)|\epsilon_{i+1}(\lambda), \, \psi_{i+1}(\lambda) | \psi_i(\lambda)$ 
for $1 \leq i \leq r-1$, also $\psi_1(\lambda) = d(\lambda)$ 
(see \cite{Rosenbrock} or \cite{Amparan-Dopico} for details). The matrix rational function 
on the right hand side of Equation \eqref{Eqn-Smith-McMillan form} is called the 
Smith-McMillan form of $R(\lambda)$. A scalar $\lambda_0 \in \mathbb{C}$ is said to be a 
finite pole of $R(\lambda)$ if $\lambda_0$ is a zero of $\psi_1(\lambda)$. If $\lambda_0$ 
is a finite pole of $R(\lambda)$, we can write 
$\psi_i(\lambda) = (\lambda - \lambda_0)^{\alpha_i} \tilde{\psi_i}(\lambda)$, 
where $\alpha_i$ is a nonnegative integer and $\tilde{\psi_i(\lambda_0)} \neq 0$ 
for $1\leq i \leq r$. The nonzero numbers in $(\alpha_1,\alpha_2,\cdots,\alpha_r)$ are 
called the partial multiplicities of $\lambda_0$ as a pole of $R(\lambda)$.  
We call the maximum partial multiplicity of $\lambda_0$ to be the order of $\lambda_0$ as 
a pole of $R(\lambda)$. In other words $\lambda_0$ is a pole of $R(\lambda)$ 
of order $k$ if it is a zero of $\psi_1(\lambda)$ of order $k$.

Recall that a scalar rational function $r(\lambda)$ can 
be expressed in the form, 
\begin{equation}\label{Eqn-scalar rational function}
\displaystyle r(\lambda) = p(\lambda) +  \frac{s(\lambda)}{q(\lambda)} 	
\end{equation}
where $p(\lambda)$, $s(\lambda)$ and $q(\lambda)$ are polynomials of degree $m$, $d_1$ 
and $d_2$ respectively with $d_1<d_2$, $q(\lambda) \not \equiv 0$ and $s(\lambda)$, $q(\lambda)$ 
are coprime. Let $q(\lambda) = c(\lambda-a_1)^{m_1} (\lambda-a_2)^{m_2} \cdots (\lambda-a_n)^{m_n}$ 
be the factorization of $q(\lambda)$ in $\mathbb{C}[\lambda]$. Then it is known that 
$\displaystyle \frac{s(\lambda)}{q(\lambda)}$ can be expressed as a polynomial in 
$\displaystyle \frac{1}{\lambda-a_j}$ for $j =1, \cdots, n$ \cite{Ponnusamy}. Therefore 
$r(\lambda)$ can be written as,
\begin{equation}\label{Eqn-scalar rational function-general form}
r(\lambda) = p(\lambda) + \displaystyle \sum_{j=1}^{n} \sum_{k=1}^{m_j} 
\frac{b_k^{(j)}}{(\lambda-a_j)^k}
\end{equation}
where, $a_j$'s denote the poles of $r(\lambda)$ with order $m_j \geq 1$ for 
$1 \leq j \leq n$. Thus, one can express an arbitrary $p \times p$ matrix rational function $R(\lambda)$ as 
\begin{equation}\label{Eqn- rational matrix}
R(\lambda) = P(\lambda) +\displaystyle \sum_{j=1}^{n} \sum_{k=1}^{m_j} 
\frac{B_k^{(j)}}{(\lambda-a_j)^k},
\end{equation}
where $P(\lambda)=C_m\lambda^m- C_{m-1}\lambda^{m-1} - \cdots -C_1 \lambda - C_0$, with 
$C_i$'s, $B_k^{(j)}$'s being square matrices of size $p\times p$ and each $a_j \in \mathbb{C}$ 
is a finite pole of $R(\lambda)$ of order $m_j$ for $1 \leq j \leq n$. \emph{Throughout 
this paper we assume that the degree of $P(\lambda)$ is at least one and that the leading coefficient of $P(\lambda)$ is the identity matrix of size $p$, that is $C_m=I_p$}. 
The particular case when $P(\lambda)$ is a linear matrix polynomial 
and each pole is of order $1$ reduces to the following form studied in \cite{Pallavi} in the context of eigenvalue location  
\begin{equation}\label{Eqn-rational matrix-linear}
R(\lambda) = \displaystyle -B_0 +I\lambda +\frac{B_1}{\lambda-a_1}+ \frac{B_2}{\lambda-a_2} + 
\cdots + \frac{B_n} {\lambda-a_n},
\end{equation} 
where $a_j$'s are distinct complex numbers. Given any $p \times p$ matrix rational 
function $R(\lambda)$ it can be expressed as $R(\lambda) = D(\lambda)+C(\lambda)A(\lambda)^{-1}B(\lambda)$, 
where $A(\lambda)$, $B(\lambda)$, $C(\lambda)$ and $D(\lambda)$ are matrix polynomials 
of size $l \times l$, $l \times p$, $p \times l$ and $p \times p$ respectively. Using 
Rosenbrock's theory (see \cite{Dopico} or \cite{Rosenbrock} for details), we get a matrix polynomial $P(\lambda)$ corresponding 
to $R(\lambda)$ of the form $P(\lambda) = \begin{bmatrix}
A(\lambda) & B(\lambda) \\
-C(\lambda) & D(\lambda)
\end{bmatrix}$, which is called a polynomial system matrix of $R(\lambda)$.  
For $R(\lambda)$ given in \eqref{Eqn-rational matrix-linear}, it is easy to see that 
the following linearization is a polynomial system matrix: 
\begin{equation}
P(\lambda) = \begin{bmatrix}
\begin{array}{c c c|c} \displaystyle
(\lambda- a_1)I & & & I \\
& \ddots & & \vdots \\
&  & (\lambda - a_m)I & I \\
\hline
-B_1 & \cdots & -B_m & I \lambda -B_0
\end{array}
\end{bmatrix}=I\lambda-C_R.
\end{equation}
This was used in \cite{Pallavi} for eigenvalue location of $R(\lambda)$ in terms of that 
of the block matrix $C_R$. However, if $P(\lambda)$ in Equation \eqref{Eqn- rational matrix} is not linear, 
then it is difficult to get a polynomial system matrix of 
$R(\lambda)$ which is in linearized form. Therefore in Section \ref{sec-3} we directly 
associate a block matrix $\mathcal{C}_R$ corresponding to $R(\lambda)$ and prove that 
the set of eigenvalues of $R(\lambda)$ is contained in the set of eigenvalues of 
$\mathcal{C}_R$. A similar technique was used by Su and Bai in \cite{Su-Bai}. 

As mentioned in the introduction, we make use of bounds on the numerical radius of a matrix 
to estimate bounds on the eigenvalues of $R(\lambda)$. For $A \in M_p$, the set 
$W(A) = \{ x^*Ax : x \in \mathbb{C}^p$ and  $x^*x=1\}$ is called the numerical range of 
$A$ and $w(A) = \sup \{ |\lambda| : \lambda \in W(A) \}$ is called the numerical radius 
of $A$. If $\rho(A)$ denotes the spectral radius of $A$, then for any eigenvalue $\lambda_0$ 
of $A$, 
\begin{equation}\label{numerical radius-relationship}
|\lambda_0| \leq \rho(A) \leq w(A) \leq ||A||_2.
\end{equation}
Other well known properties of the numerical radius are  
\begin{itemize}
\item $w(A) \geq 0$.
\item $w(\alpha A) = |\alpha| w(A)$ for every $\alpha \in \mathbb{C}$.
\item $w(A+B) \leq w(A)+w(B)$ for $A$, $B$ in $M_p$. 
\end{itemize}

Before stating our main results we state without proof bounds on the numerical radius of matrices.

\begin{lemma}[\cite{Abu-Kittaneh}] \label{Lem-numerical radius-1}
Let $A \in M_k$, $B \in M_{k,s}$, $C \in M_{s,k}$ and $D \in M_{s}$, and 
let $K = \begin{bmatrix}
A & B \\
C & D
\end{bmatrix}$. Then 
$w(K) \leq \displaystyle \frac{1}{2} \Big(w(A)+ w(D) +\sqrt{(w(A)-w(D))^2+ 4w^2(K_0)}\Big)$,
where $K_0 = \begin{bmatrix}
0 & B \\
C & 0
\end{bmatrix}$. 
\end{lemma}

\begin{lemma}[\cite{Abu-Kittaneh}]\label{Lem-numerical radius-2}
Let $B \in M_{k,s}$, $C \in M_{s,k}$, then
$\displaystyle w\Bigg(\begin{bmatrix}
0 & B \\
C & 0
\end{bmatrix}\Bigg) \leq \frac{1}{2}\big(||B||_2+ ||C||_2 \big)$.
\end{lemma}

\medskip
\noindent
The following Corollary is a consequence of Lemmas \ref{Lem-numerical radius-1} and \ref{Lem-numerical radius-2}.

\begin{corollary}\label{Cor-numerical radius}
Let $A \in M_k$, $B \in M_{k,s}$, $C \in M_{s,k}$ and $D \in M_{s}$, and 
let $K = \begin{bmatrix}
A & B \\
C & D
\end{bmatrix}$. Then $w(K) \leq \displaystyle \frac{1}{2} \Big(w(A)+ w(D) +\sqrt{(w(A)-w(D))^2+ 
(||B||_2+||C||_2)^2} \Big)$.
\end{corollary}

\medskip
\begin{lemma} [\cite{R Bhatia}]\label{Lem-numerical radius-3}
Let $T = \text{diag}(T_{11}, T_{22}, \ldots, T_{nn})$, where $T_{ii} \in M_{n_i}$ for 
$1 \leq i \leq n$. Then $w(T) =\displaystyle \max_{1 \leq i \leq n}\{w(T_{ii})\}$.
\end{lemma}

It is also not hard to prove the following result as in the example given in Section $1.3$ of 
\cite{Gustafson-Rao}. We state this as a lemma for ready reference.

\begin{lemma}\label{Lem-numerical radius-Jordan block}
Let $L_n$ be $n \times n$ matrix given by $L_n =\begin{bmatrix}
a & 0 & \cdots & 0 & 0\\
1 & a & \cdots & 0 & 0\\
\vdots & \vdots & \ddots & \vdots & \vdots \\
0 & 0 & \cdots & 1 & a 
\end{bmatrix}$, where $a$ is a complex number. Then $w(L_n) \leq \displaystyle |a|+ 
\cos\Big(\frac{\pi}{n+1}\Big)$. 
If $a$ is a nonnegative real number then equality holds.
\end{lemma}

\medskip
\section{Main results}\label{sec-3}
This section contains the main results of this manuscript. For ease of reading the 
section is further divided into subsections.

\medskip
\subsection{Associated block matrix corresponding to matrix rational function}\label{sec-3.1}

Consider the matrix rational function $R(\lambda)$ given in Equation 
\eqref{Eqn- rational matrix}. Corresponding to $R(\lambda)$ we associate a block 
matrix $\mathcal{C}_R$ of size $p\Big(m + \displaystyle \sum_{j=1}^{n}\frac{m_j(m_j+1)}{2}\Big)$ 
defined as follows.
\begin{equation}
\mathcal{C}_R= \begin{bmatrix}
\mathcal{A}_n & 0 & \cdots & 0 & -\mathcal{F}_n \\
0 & \mathcal{A}_{n-1} & \cdots & 0 & -\mathcal{F}_{n-1} \\
\vdots & \vdots & \ddots & \vdots & \vdots \\
0 & 0 & \cdots & \mathcal{A}_1 & -\mathcal{F}_1 \\
\mathcal{B}_n & \mathcal{B}_{n-1} & \cdots & \mathcal{B}_1 & \mathcal{B}_0
\end{bmatrix},
\end{equation}
where for each $1 \leq j \leq n$ the blocks are defined as follows: \\
$\mathcal{A}_j = \text{diag}(\textbf{A}_{m_j}^{(j)},\textbf{A}_{m_j-1}^{(j)},
\cdots,\textbf{A}_1^{(j)})$ is a block diagonal matrix of size $ \displaystyle 
\frac{pm_j(m_j+1)}{2}$, with $\textbf{A}_k^{(j)} 
=\begin{bmatrix}
a_jI & I & 0 & \cdots & 0 & 0 \\
0 & a_jI & I & \cdots & 0 & 0 \\
\vdots & \vdots & \vdots & \ddots & \vdots & \vdots \\
0 & 0 & 0 & \cdots & a_jI & I \\
0 & 0 & 0 & \cdots & 0 & a_jI
\end{bmatrix}_{pk \times pk}$, $1\leq k \leq m_j$ and 

\noindent
$\mathcal{B}_j = \begin{bmatrix}
\textbf{B}_{m_j}^{(j)} & \textbf{B}_{m_j-1}^{(j)} & \cdots & \textbf{B}_1^{(j)}
\end{bmatrix}$ is a block row matrix of size $\displaystyle pm \times 
\frac{pm_j(m_j+1)}{2}$, with $\textbf{B}_k^{(j)} = \begin{bmatrix}
0 & 0 & \cdots & 0 \\
0 & 0 & \cdots & 0 \\
\vdots & \vdots & \ddots & \vdots \\
B_{k}^{(j)} & 0 & \cdots & 0
\end{bmatrix}_{pm \times pk}, 1 \leq k \leq m_j $, \\
$\mathcal{B}_0 = \begin{bmatrix}
0 & I & 0 & \cdots & 0\\
0 & 0 & I & \cdots & 0\\
\vdots & \vdots & \vdots & \ddots & \vdots\\
0 & 0 & 0 & \cdots & I\\
C_0 & C_1 & C_2 & \cdots & C_{m-1}\\ 
\end{bmatrix}_{pm \times pm}$ and
$\mathcal{F}_j = \begin{bmatrix}
\textbf{F}_{m_j} \\
\textbf{F}_{m_j-1} \\
\vdots \\
\textbf{F}_1
\end{bmatrix}$ is block column matrix of size $\displaystyle \frac{pm_j(m_j+1)}{2} 
\times pm$ with $ \textbf{F}_k =\begin{bmatrix}
0 & 0 & \cdots & 0\\
0 & 0 & \cdots & 0\\
\vdots & \vdots & \ddots & \vdots \\
I & 0 & \cdots & 0 \\
\end{bmatrix}_{pk \times pm}, 1 \leq k \leq m_j$ for $j = 1,\cdots, n$.

\medskip
Our first result is the following: the set of eigenvalues of $R(\lambda)$ is contained 
in the set of eigenvalues of $\mathcal{C}_R$.

\begin{theorem}\label{Thm- eigenvalues- matrix rational function}
Let $R(\lambda)$ be the matrix rational function given in Equation \eqref{Eqn- rational matrix} 
and $\mathcal{C}_R$ be the associated block 
matrix described above. Then the eigenvalues of $R(\lambda)$ are also eigenvalues of 
$\mathcal{C}_R$.
\end{theorem}

\begin{proof}
Let $(\lambda_0,v)$ be an eigenpair of $R(\lambda)$ so that 
$R(\lambda_0)v = 0$.  
For $1 \leq j \leq n$, define a vector $\mathcal{Y}_j=\begin{bmatrix}
\textbf{Y}_{m_j}^{(j)} \\
\textbf{Y}_{m_j -1}^{(j)} \\
\vdots \\
\textbf{Y}_{2}^{(j)} \\
\textbf{Y}_{1}^{(j)}\\
\end{bmatrix}$, where $\textbf{Y}_k^{(j)} = \displaystyle \begin{bmatrix}
\frac{-1}{(\lambda_0-a_j)^k} v\\
\frac{-1}{(\lambda_0-a_j)^{k-1}} v\\
\vdots\\
\frac{-1}{(\lambda_0-a_j)^2}v\\
\frac{-1}{(\lambda_0-a_j)}v
\end{bmatrix}$ for $1 \leq k \leq m_j$ and $\mathcal{Y}_0 = \begin{bmatrix}
v\\
\lambda_0v\\
\vdots\\
\lambda_0^{m-1}v\\
\end{bmatrix}$. Let $\mathcal{Y} =\begin{bmatrix}
\mathcal{Y}_n \\
\mathcal{Y}_{n-1} \\
\vdots \\
\mathcal{Y}_1 \\
\mathcal{Y}_0\\
\end{bmatrix}$. Consider 
\begin{align*}
\mathcal{C}_R \mathcal{Y} = &\begin{bmatrix}
\mathcal{A}_n & 0 & \cdots & 0 & -\mathcal{F}_n \\
0 & \mathcal{A}_{n-1} & \cdots & 0 & -\mathcal{F}_{n-1} \\
\vdots & \vdots & \ddots & \vdots & \vdots \\
0 & 0 & \cdots & \mathcal{A}_1 & -\mathcal{F}_1 \\
\mathcal{B}_n & \mathcal{B}_{n-1} & \cdots & \mathcal{B}_1 & \mathcal{B}_0
\end{bmatrix} \begin{bmatrix}
\mathcal{Y}_n \\
\mathcal{Y}_{n-1} \\
\vdots \\
\mathcal{Y}_1 \\
\mathcal{Y}_0\\
\end{bmatrix} 
= \begin{bmatrix} 
\mathcal{A}_n \mathcal{Y}_n-\mathcal{F}_n\mathcal{Y}_0 \\
\mathcal{A}_{n-1} \mathcal{Y}_{n-1}-\mathcal{F}_{n-1}\mathcal{Y}_0 \\
\vdots \\
\mathcal{A}_1 \mathcal{Y}_1-\mathcal{F}_1\mathcal{Y}_0 \\
\displaystyle \sum_{j=1}^{n} \mathcal{B}_j \mathcal{Y}_j + \mathcal{B}_0\mathcal{Y}_0\\
\end{bmatrix}.
\end{align*}
For $1 \leq j \leq n$ we have, 
\begin{align*}
\mathcal{A}_{j} \mathcal{Y}_j-\mathcal{F}_{j}\mathcal{Y}_0= &\begin{bmatrix}
\textbf{A}_{m_j}^{(j)} & 0 & \cdots & 0 & 0 \\
0 & \textbf{A}_{m_j -1}^{(j)} & \cdots & 0 & 0 \\
\vdots & \vdots & \ddots & \vdots & \vdots \\
0 & 0 & \cdots & \textbf{A}_{2}^{(j)} & 0 \\
0 & 0 & \cdots & 0 & \textbf{A}_{1}^{(j)} \\
\end{bmatrix} \begin{bmatrix}
\textbf{Y}_{m_j}^{(j)} \\
\textbf{Y}_{m_j -1}^{(j)} \\
\vdots \\
\textbf{Y}_{2}^{(j)} \\
\textbf{Y}_{1}^{(j)}\\
\end{bmatrix}- \begin{bmatrix}
\textbf{F}_{m_j} \\
\textbf{F}_{m_j-1} \\
\vdots \\
\textbf{F}_2\\
\textbf{F}_1
\end{bmatrix}\mathcal{Y}_0\\ 
=& \begin{bmatrix}
\textbf{A}_{m_j}^{(j)}\textbf{Y}_{m_j}^{(j)}-\textbf{F}_{m_j}\mathcal{Y}_0\\
\textbf{A}_{m_j-1}^{(j)}\textbf{Y}_{m_j-1}^{(j)}-\textbf{F}_{m_j-1}\mathcal{Y}_0\\
\vdots\\
\textbf{A}_2^{(j)}\textbf{Y}_2^{(j)}-\textbf{F}_2\mathcal{Y}_0\\
\textbf{A}_1^{(j)}\textbf{Y}_1^{(j)}-\textbf{F}_1\mathcal{Y}_0
\end{bmatrix}.
\end{align*} 
For $1 \leq k \leq m_j$ consider  
\begin{align*}
& \textbf{A}_k^{(j)}\textbf{Y}_k^{(j)}-\textbf{F}_k\mathcal{Y}_0 \nonumber\\
& = \begin{bmatrix}
a_jI & I & \cdots & 0 & 0 \\
0 & a_jI  & \cdots & 0 & 0 \\
\vdots & \vdots & \ddots & \vdots & \vdots \\
0 & 0 & \cdots & a_jI & I \\
0 & 0 & \cdots & 0 & a_jI
\end{bmatrix} \begin{bmatrix}
\frac{-1}{(\lambda_0-a_j)^k} v\\
\frac{-1}{(\lambda_0-a_j)^{k-1}} v\\
\vdots\\
\frac{-1}{(\lambda_0-a_j)^2}v\\
\frac{-1}{(\lambda_0-a_j)}v
\end{bmatrix}-\begin{bmatrix}
0 & 0 & \cdots & 0\\
0 & 0 & \cdots & 0\\
\vdots & \vdots & \ddots & \vdots \\
0 & 0 & \cdots &  0 \\
I & 0 & \cdots & 0 \\
\end{bmatrix} \begin{bmatrix}
v\\
\lambda_0v\\
\vdots\\
\lambda_0^{m-2}v \\
\lambda_0^{m-1}v\\
\end{bmatrix} \nonumber\\ 
& = \lambda_0 \begin{bmatrix}
\frac{-1}{(\lambda_0-a_j)^k} v\\
\frac{-1}{(\lambda_0-a_j)^{k-1}} v\\
\vdots\\
\frac{-1}{(\lambda_0-a_j)^2}v\\
\frac{-1}{(\lambda_0-a_j)}v
\end{bmatrix} = \lambda_0 \textbf{Y}_k^{(j)}. 
\end{align*}
This implies 
\begin{equation}
\mathcal{A}_{j} \mathcal{Y}_j-\mathcal{F}_{j}\mathcal{Y}_0  =\lambda_0 
\mathcal{Y}_j  \ \text{for} \ j= 1, 2, \cdots, n. 
\end{equation}

Now consider,
\begin{align*}
& \displaystyle \sum_{j=1}^{n} \mathcal{B}_j \mathcal{Y}_j+ \mathcal{B}_0\mathcal{Y}_0 \\ 
& = \sum_{j=1}^{n} \begin{bmatrix}
\textbf{B}_{m_j}^{(j)} & \textbf{B}_{m_j -1}^{(j)} & \cdots & \textbf{B}_1^{(j)}  
\end{bmatrix} \begin{bmatrix}
\textbf{Y}_{m_j}^{(j)}\\
\textbf{Y}_{m_j -1}^{(j)}\\
\vdots \\
\textbf{Y}_1^{(j)}\\
\end{bmatrix}+\mathcal{B}_0\mathcal{Y}_0 \\
%& = \sum_{j=1}^{n} \sum_{k=1}^{m_j} \textbf{B}_k^{(j)}\textbf{Y}_k^{(j)} +
%\mathcal{B}_0\mathcal{Y}_0 \\
& =\mathlarger \sum_{j=1}^{n} \mathlarger \sum_{k=1}^{m_j} \begin{bmatrix}
0 & 0 & \cdots & 0 \\
0 & 0 & \cdots & 0 \\
\vdots & \vdots & \ddots & \vdots\\
B_k^{(j)} & 0 & \cdots & 0 \\
\end{bmatrix} \begin{bmatrix}
\frac{-1}{(\lambda_0-a_j)^k}v\\
\frac{-1}{(\lambda_0-a_j)^{k-1}}v\\
\vdots\\
\frac{-1}{(\lambda_0-a_j)}v\\
\end{bmatrix}
\end{align*} 

$+ \begin{bmatrix}
0 & I & 0 & \cdots & 0\\
0 & 0 & I & \cdots & 0\\
\vdots & \vdots & \vdots & \ddots & \vdots\\
C_0 & C_1 & C_2 & \cdots & C_{m-1}  
\end{bmatrix}\begin{bmatrix}
v\\
\lambda_0v\\
\vdots\\
\lambda_0^{m-1}v
\end{bmatrix}$ \\
%\end{align*}
\begin{align}
& = \sum_{j=1}^{n} \sum_{k=1}^{m_j} \begin{bmatrix}
0\\
0\\
\vdots\\
\frac{-B_k^{(j)}v}{(\lambda_0-a_j)^k}\\
\end{bmatrix}+\begin{bmatrix}
\lambda_0v\\
\lambda_0^2v\\
\vdots\\
C_0v+C_1\lambda_0v+ \cdots + C_{m-1}\lambda_0^{m-1}v
\end{bmatrix} \nonumber\\ 
%& =\begin{bmatrix}
%\lambda_0v\\
%\lambda_0^2v\\
%\vdots\\
%\displaystyle \mathlarger \sum_{j=1}^{n} \mathlarger \sum_{k=1}^{m_j}\frac{-B_k^{(j)}}{(\lambda_0-a_j)^k}
%+C_0v+C_1\lambda_0v+ \cdots + C_{m-1}\lambda_0^{m-1}v
%\end{bmatrix}  \nonumber\\
& =\begin{bmatrix}
\lambda_0v\\
\lambda_0^2v\\
\vdots\\
\lambda_0^mv\\
\end{bmatrix}
=\lambda_0 \begin{bmatrix}
v\\
\lambda_0v\\
\vdots\\
\lambda_0^{m-1}v\\
\end{bmatrix}=\lambda_0\mathcal{Y}_0.
\end{align}
Therefore from $(3.2)$ and $(3.3)$ we get $\mathcal{C}_R \mathcal{Y} 
= \lambda_0 \mathcal{Y}$. Thus, if $\lambda_0$ is an eigenvalue of $R(\lambda)$, 
it is an eigenvalue of $\mathcal{C}_R$ also.
\end{proof}

\subsection{Bounds on the zeros of scalar rational functions}\label{sec-3.2}

In this section we consider a scalar rational function $r(\lambda)$ and derive  
bounds on its zeros by associating a suitable block matrix. Consider the scalar rational function 
\begin{equation}\label{Eqn-scalar rational function}
r(\lambda) = \lambda^m- c_{m-1}\lambda^{m-1} - \cdots -c_1 \lambda - c_0 -  
\displaystyle \sum_{j=1}^{n} \sum_{k=1}^{m_j} \frac{b_k^{(j)}}{(\lambda-a_j)^k}, 
\end{equation}
where $b_k^{(j)}$'s, $a_j$'s and $c_i$'s are complex numbers.
Since $r(\lambda)$ can be considered as a matrix rational function, corresponding 
to $r(\lambda)$ we have the associated block matrix of size $m + \displaystyle 
\sum_{j=1}^{n}\frac{m_j(m_j+1)}{2}$ 
\begin{center}
$\mathcal{C}_r= \begin{bmatrix}
\mathcal{A}_n & 0 & \cdots & 0 & -\mathcal{F}_n \\
0 & \mathcal{A}_{n-1} & \cdots & 0 & -\mathcal{F}_{n-1} \\
\vdots & \vdots & \ddots & \vdots & \vdots \\
0 & 0 & \cdots & \mathcal{A}_1 & -\mathcal{F}_1 \\
\mathcal{B}_n & \mathcal{B}_{n-1} & \cdots & \mathcal{B}_1 & \mathcal{B}_0
\end{bmatrix}$, 
\end{center}
where for each $1 \leq j \leq n$,   
$\mathcal{A}_j = \text{diag}(A_{m_j}^{(j)},A_{m_j-1}^{(j)},\cdots,A_1^{(j)})$ is a 
block diagonal matrix of size $\displaystyle \frac{m_j(m_j+1)}{2}$ with $A_k^{(j)} 
=\begin{bmatrix}
a_j & 1 & 0 & \cdots & 0 & 0 \\
0 & a_j & 1 & \cdots & 0 & 0 \\
\vdots & \vdots & \vdots & \ddots & \vdots & \vdots \\
0 & 0 & 0 & \cdots & a_j & 1 \\
0 & 0 & 0 & \cdots & 0 & a_j
\end{bmatrix}_{k \times k}$ for $1\leq k \leq m_j$,
$\mathcal{B}_j = \begin{bmatrix}
B_{m_j}^{(j)} & B_{m_j-1}^{(j)} & \cdots & B_1^{(j)}
\end{bmatrix}$ is a block row matrix of size $\displaystyle m \times 
\frac{m_j(m_j+1)}{2}$ with $B_k^{(j)} = \begin{bmatrix}
0 & 0 & \cdots & 0 \\
0 & 0 & \cdots & 0 \\
\vdots & \vdots & \ddots & \vdots \\
-b_{k}^{(j)} & 0 & \cdots & 0
\end{bmatrix}_{m \times k}$ for $ 1 \leq k \leq m_j $, 
$\mathcal{B}_0 = \begin{bmatrix}
0 & 1 & 0 & \cdots & 0\\
0 & 0 & 1 & \cdots & 0\\
\vdots & \vdots & \vdots & \ddots & \vdots\\
0 & 0 & 0 & \cdots & 1\\
c_0 & c_1 & c_2 & \cdots & c_{m-1}\\ 
\end{bmatrix}_{m \times m}$, and
$\mathcal{F}_j = \begin{bmatrix}
F_{m_j} \\
F_{m_j-1} \\
\vdots \\
F_1
\end{bmatrix}$ is a block column matrix of size $\displaystyle \frac{m_j(m_j+1)}{2} 
\times m$ with $ F_k =\begin{bmatrix}
0 & 0 & \cdots & 0\\
0 & 0 & \cdots & 0\\
\vdots & \vdots & \ddots & \vdots \\
1 & 0 & \cdots & 0 \\
\end{bmatrix}_{k \times m}, 1 \leq k \leq m_j$ for $j = 1,\cdots, n$.

\medskip
\noindent
The following corollary to Theorem \ref{Thm- eigenvalues- matrix rational function} is 
immediate.

\begin{corollary}\label{Cor-zeros of rational function}
Let $r(\lambda)$ be a rational function as in Equation \eqref{Eqn-scalar rational function} and let $\mathcal{C}_r$ 
be the associated block matrix. Then the zeros of $r(\lambda)$ are eigenvalues of $\mathcal{C}_r$.
\end{corollary}

\medskip

We now give bounds on the zeros of $r(\lambda)$ using the norm of the associated block 
matrix.

\begin{theorem}\label{Thm-bounds for zeros using norms}
Let $r(\lambda)$ be a rational function as in Equation \eqref{Eqn-scalar rational function}. If 
$\lambda_0$ is a zero of $r(\lambda)$ then,
\begin{enumerate}
\item $\displaystyle  |\lambda_0| \leq \max_{1 \leq j \leq n} \bigg\{1+|a_j|, 
\sum_{j=1}^{n} \sum_{k=1}^{m_j}|b_k^{(j)}|+ \sum_{i=0}^{m-1}|c_i| \bigg\}$.

\item 
$ |\lambda_0| \leq 
\displaystyle \max_{\substack{\text{$1 \leq j \leq n$}\\ \text {$1 \leq i \leq m-1$}}} 
\bigg\{|a_j|+|b^{(j)}|, 1+|c_i|, |c_0|+ n \bigg\}$,  if each pole is of order $1$.\\


\ \ $\leq \displaystyle \max_{\substack{\text{$1 \leq j \leq n$}\\ 
\text{$1 \leq k \leq m_j$} \\ \text {$1 \leq i \leq m-1$}}} 
\bigg\{|a_j|+|b_k^{(j)}|, 1+|a_j|, 1+|c_i|, |c_0|+ \sum_{j=1}^{n}m_j \bigg\}$, otherwise.

%\item $\displaystyle |\lambda_0| \leq 
%\max_{\substack{1 \leq j \leq n \\ 1 \leq i \leq m-1}} \bigg\{|a_j|+|b^{(j)}|, 1+|c_i|, 
%|c_0|+ n \bigg\}, \ \text{if each $a_j$ is a pole of order one}$.
%
%\item $\displaystyle |\lambda_0| \leq
%\max_{\substack{1 \leq j \leq n \\ 1 \leq k \leq m_j \\ 
%1 \leq i \leq m-1}} \bigg\{|a_j|+|b_k^{(j)}|, 1+|a_j|, 1+|c_i|, |c_0|+ \sum_{j=1}^{n}m_j 
%\bigg\}$, otherwise.

\end{enumerate}
\end{theorem}

\begin{proof}
Since $\lambda_0$ is a zero of $r(\lambda)$, by Corollary \ref{Cor-zeros of rational function} 
$\lambda_0$ is an eigenvalue of the corresponding block matrix $\mathcal{C}_r$. One 
can easily see that, 
\begin{align*}
||\mathcal{C}_r||_\infty \displaystyle= \max_{1 \leq j \leq n} \bigg\{1+|a_j|, 
\sum_{j=1}^{n} \sum_{k=1}^{m_j}|b_k^{(j)}|+ \sum_{i=0}^{m-1}|c_i| \bigg\}.
\end{align*}
We also have 
$||\mathcal{C}_r||_1 \displaystyle = \max_{\substack{1 \leq j \leq n \\ 1 \leq i \leq m-1}} 
\bigg\{|a_j|+|b^{(j)}|, 1+|c_i|, |c_0|+ n 
\bigg\}$, if each $a_j$ is a pole of order one and equals 
$\displaystyle \max_{\substack{1 \leq j \leq n \\ 1 \leq k \leq m_j\\ 1 \leq i 
\leq m-1}} \bigg\{|a_j|+|b_k^{(j)}|, 1+|a_j|, 1+|c_i|, |c_0|+ \sum_{j=1}^{n}m_j \bigg\}$ 
otherwise. Hence the proof.
\end{proof}

In the two theorems that follow, we use bounds on the numerical radius of the associated block matrix to 
give bounds on the zeros of $r(\lambda)$.

\begin{theorem}\label{Thm-bounds on zeros using numerical range-2}
Let $r(\lambda)$ be a rational function as in Equation \eqref{Eqn-scalar rational function}, 
where the $a_j$'s are nonnegative reals for $1 \leq j \leq n$. 
If $\lambda_0$ is a zero of $r(\lambda)$, then 
\begin{align*}
|\lambda_0| \leq \displaystyle \frac{\alpha+ \beta +\sqrt{(\alpha-\beta)^2 + (\gamma+\delta)^2}}{2},
\end{align*}
where $\alpha = \displaystyle \max_{1 \leq j \leq n} \Big\{a_j + 
\cos\Big(\frac{\pi}{m_j+1}\Big) \Big\}$, $\beta = w(\mathcal{B}_0)$, $\gamma = 
\displaystyle \sqrt{\sum_{j=1}^{n} m_j}$ and $\delta = \displaystyle 
\sqrt{ \sum_{j=1}^{n}\sum_{k=1}^{m_j} {|b_{k}^{(j)}|}^2}$ with $\mathcal{B}_0 = 
\begin{bmatrix}
0 & 1 & 0 & \cdots & 0\\
0 & 0 & 1 & \cdots & 0\\
\vdots & \vdots & \vdots & \ddots & \vdots\\
0 & 0 & 0 & \cdots & 1\\
c_0 & c_1 & c_2 & \cdots & c_{m-1}\\ 
\end{bmatrix}_{m \times m}$.
\end{theorem}

\begin{proof}
Since $\lambda_0$ is a zero of $r(\lambda)$, by Corollary \ref{Cor-zeros of rational function} 
$\lambda_0$ is an eigenvalue of the corresponding block matrix $\mathcal{C}_r= \begin{bmatrix}
\begin{array}{c c c c|c}
\mathcal{A}_n & 0 & \cdots & 0 & -\mathcal{F}_n \\
0 & \mathcal{A}_{n-1} & \cdots & 0 & -\mathcal{F}_{n-1} \\
\vdots & \vdots & \ddots & \vdots & \vdots \\
0 & 0 & \cdots & \mathcal{A}_1 & -\mathcal{F}_1 \\
\hline
\mathcal{B}_n & \mathcal{B}_{n-1} & \cdots & \mathcal{B}_1 & \mathcal{B}_0
\end{array}
\end{bmatrix} = \begin{bmatrix}
A & C \\
B & D \\
\end{bmatrix}$. By Corollary \ref{Cor-numerical radius} we have, 
\begin{align*}
w(\mathcal{C}_r) & \leq \frac{w(A)+ w(D) +\sqrt{(w(A)-w(D))^2 + (||C||_2+||B||_2)^2}}{2}.
\end{align*} 
We also have from Lemma \ref{Lem-numerical radius-Jordan block} that 
$w(A)= \displaystyle \max_{1 \leq j \leq n} 
\Big\{a_j + \cos\Big(\frac{\pi}{m_j+1}\Big) \Big\}$. 
Note that for $1 \leq j \leq n$ and $1 \leq k \leq m_j$, $B_k^{(j)}{B_k^{(j)}}^* 
= \begin{bmatrix}
0 & 0 & \cdots & 0 \\
0 & 0 & \cdots & 0 \\
\vdots & \vdots & \ddots & \vdots \\
0 & 0 & \cdots & {|b_{k}^{(j)}|}^2
\end{bmatrix}_{m \times m}$. 
Therefore $\mathcal{B}_j\mathcal{B}_j^* = \begin{bmatrix}
B_{m_j}^{(j)} & B_{m_j-1}^{(j)} & \cdots & B_1^{(j)}
\end{bmatrix}\begin{bmatrix}
{B_{m_j}^{(j)}}^* \\ 
{B_{m_j-1}^{(j)}}^* \\
\vdots \\ 
{B_1^{(j)}}^*
\end{bmatrix} = B_{m_j}^{(j)}{B_{m_j}^{(j)}}^*+ B_{m_j-1}^{(j)}{B_{m_j-1}^{(j)}}^* 
+ \cdots + B_1^{(j)}{B_1^{(j)}}^* = \begin{bmatrix}
0 & 0 & \cdots & 0 \\
0 & 0 & \cdots & 0 \\
\vdots & \vdots & \ddots & \vdots \\
0 & 0 & \cdots & \displaystyle \sum_{k=1}^{m_j} {|b_{k}^{(j)}|}^2
\end{bmatrix}_{m \times m}$. One can then compute $B B^* = \begin{bmatrix}
\mathcal{B}_n & \mathcal{B}_{n-1} & \cdots & \mathcal{B}_1
\end{bmatrix} \begin{bmatrix}
\mathcal{B}_n^* \\ 
\mathcal{B}_{n-1}^* \\
\vdots \\
\mathcal{B}_1^*
\end{bmatrix}= \mathcal{B}_n\mathcal{B}_n^*+ \cdots + \mathcal{B}_1\mathcal{B}_1^* 
=\begin{bmatrix}
0 & 0 & \cdots & 0 \\
0 & 0 & \cdots & 0 \\
\vdots & \vdots & \ddots & \vdots \\
0 & 0 & \cdots & \displaystyle \sum_{j=1}^{n}\sum_{k=1}^{m_j} {|b_{k}^{(j)}|}^2
\end{bmatrix}$.
Thus $||B||_2 = \sqrt{\displaystyle \sum_{j=1}^{n}\sum_{k=1}^{m_j} {|b_{k}^{(j)}|}^2} $. 
Similarly, $||C||_2= \displaystyle \sqrt{\sum_{j=1}^{n}\sum_{k=1}^{m_j}1} = 
\sqrt{\sum_{j=1}^{n} m_j}$. This completes the proof.
\end{proof}

\medskip
\subsection{Bounds on the eigenvalues $R(\lambda)$ using a real rational function} 

In this section we consider a real rational function $r(\lambda)$ associated with 
$R(\lambda)$ and prove that a particular root of $r(\lambda)$ is a bound on the 
eigenvalues of the matrix rational function $R(\lambda)$. We also make use of bounds 
on the zeros of scalar rational function obtained in Section \ref{sec-3.2} to give 
bounds on the eigenvalues of $R(\lambda)$. We begin with the following lemma, the proof 
of which is a simple consequence of the intermediate value theorem.
 
\begin{lemma}\label{Lem-zeros of rational function}
Let $q(x) = x^m-\alpha_{m-1}x^{m-1}-\cdots - \alpha_1 x - \alpha_0 - \displaystyle 
\sum_{j=1}^{n} \sum_{k=1}^{m_j} \frac{\beta_k^{(j)}}{(x-\gamma_j)^k}$ be a real 
rational function, where the $\beta_k^{(j)}$'s are positive real numbers, $\alpha_i$'s 
and $\gamma_j$'s are nonnegative real numbers such that $\gamma_1 < \gamma_2 < \cdots 
< \gamma_n$. Then $q(x)$ has a root $R$ such that $ \gamma_n < R$.
\end{lemma}

%\begin{proof}
%Note that $q(x)$ is continuous in the open interval $(\gamma_n, + \infty)$ and 
%$\displaystyle \lim_{x \rightarrow \gamma_n+} = - \infty$ and $\displaystyle 
%\lim_{x \rightarrow + \infty} = + \infty$. The result then follows from 
%the intermediate value theorem.
%\end{proof}

\medskip

\begin{theorem}\label{Thm-eigenvalue bounds-matrix rational function}
Let $R(\lambda)$ be a $p \times p$ matrix rational function as in Equation \eqref{Eqn- rational matrix}. 
Associate a real rational function corresponding to $R(\lambda)$ as follows:
\begin{center}
$q(x) = x^m-||C_{m-1}||x^{m-1}-\cdots - ||C_1|| x - ||C_0|| - \displaystyle 
\sum_{j=1}^{n} \sum_{k=1}^{m_j} \frac{||B_k^{(j)}||}{(x- |a_j|)^k}$,
\end{center}
where $||\cdot ||$ is any induced matrix norm. If $\lambda_0$ is an eigenvalue of 
$R(\lambda)$ then $|\lambda_0| \leq R$, where $R$ is a real root of $q(x)$ such that 
$|a_j| < R$ for all $j=1,\cdots,n$.
\end{theorem}

\begin{proof}
Without loss of generality assume that $|a_1| < |a_2| < \cdots < |a_n|$. By Lemma 
\ref{Lem-zeros of rational function}, $q(x)$ has a real root $R$ such that $|a_j| < R$ 
for all $j=1,\cdots,n$. We 
then have 
\begin{align*}
0 & = q(R)=R^m - ||C_{m-1}|| R^{m-1} - \cdots - ||C_1|| R -||C_0||-\displaystyle 
\sum_{j=1}^{n} \sum_{k=1}^{m_j} \frac{||B_k^{(j)}||}{(R- |a_j|)^k}.
\end{align*} 
This implies $R^m  = ||C_{m-1}|| R^{m-1} + \cdots + ||C_1|| R + ||C_0||+\displaystyle 
\sum_{j=1}^{n} \sum_{k=1}^{m_j} \frac{||B_k^{(j)}||}{(R- |a_j|)^k}$. Therefore
\begin{align*}
1 = R^{-1}||C_{m-1}||  + \cdots + R^{-(m-1)} ||C_1|| + R^{-m} ||C_0||
+\displaystyle R^{-m} \sum_{j=1}^{n} \sum_{k=1}^{m_j} \frac{||B_k^{(j)}||}{(R- |a_j|)^k}.
\end{align*}
Suppose $\lambda_0 \in \mathbb{C}$ is such that $|\lambda_0| > R$. Then 
$|\lambda_0-a_j|\geq |\lambda_0|-|a_j|>R-|a_j|$. Since $|a_j| < R$, we have 
$\displaystyle \frac{1}{(|\lambda_0-a_j|)^k} < \frac{1}{(R-|a_j|)^k}$ for all 
$j=1,2,\cdots,n$ and $k=1,2,\cdots,m_j$. For any unit vector $v \in \mathbb{C}^n$ 
consider, 
\begin{align*}
& \Big|\Big| - \lambda_0^{-1}C_{m-1}v - \cdots - \lambda_0^{-(m-1)} C_1 v - 
\lambda_0^{-m} C_0 v+ \displaystyle \lambda_0^{-m} \sum_{j=1}^{n} \sum_{k=1}^{m_j} 
\frac{B_k^{(j)}}{(\lambda_0-a_j)^k} v \Big|\Big|\\
& \leq |\lambda_0|^{-1} ||C_{m-1}|| + \cdots + 
|\lambda_0|^{-m} ||C_0|| + \displaystyle |\lambda_0|^{-m} \sum_{j=1}^{n} \sum_{k=1}^{m_j} 
\frac{||B_k^{(j)}||}{(|\lambda_0-a_j|)^k} \\
& < R^{-1} ||C_{m-1}|| + \cdots + R^{-(m-1)} ||C_1|| + R^{-m} ||C_0|| + \displaystyle 
R^{-m} \sum_{j=1}^{n} \sum_{k=1}^{m_j} \frac{||B_k^{(j)}||}{(R-|a_j|)^k} \\
& =1.
\end{align*}
Now consider 
\begin{align*}
& ||R(\lambda_0)v|| \\
& =\Big|\Big| \lambda_0^m v- C_{m-1}\lambda_0^{m-1}v - \cdots -C_1 \lambda_0 v - 
C_0 v+ \displaystyle \sum_{j=1}^{n} \sum_{k=1}^{m_j} 
\frac{B_k^{(j)}}{(\lambda_0-a_j)^k} v\Big|\Big| \\
& = |\lambda_0|^m\Big|\Big|v- \lambda_0^{-1}C_{m-1}v - \cdots - 
\lambda_0^{-m} C_0 v+ \displaystyle \lambda_0^{-m} 
\sum_{j=1}^{n} \sum_{k=1}^{m_j} \frac{B_k^{(j)}}{(\lambda_0-a_j)^k} v \Big|\Big| \\
& \geq |\lambda_0|^m\Big(1-\Big|\Big| - \lambda_0^{-1}C_{m-1}v - \cdots - 
\lambda_0^{-m} C_0 v+ \displaystyle \lambda_0^{-m} 
\sum_{j=1}^{n} \sum_{k=1}^{m_j} \frac{B_k^{(j)}}{(\lambda_0-a_j)^k} v \Big|\Big| 
\Big)\\
& > |\lambda_0|^m(1-1) = 0.
\end{align*}
Therefore if $\lambda_0\in \mathbb{C}$ is an eigenvalue of $R(\lambda)$, then 
$|\lambda_0| \leq R$. This completes the proof.
\end{proof}

\medskip
The following theorem summarizes what has been proved so far.

\begin{theorem}\label{Thm summary}
Let $R(\lambda)$ be a $p \times p$ matrix rational function as in Equation \eqref{Eqn- rational matrix}. 
If $\lambda_0$ be an eigenvalue of $R(\lambda)$, then,

\begin{enumerate}
\item $\displaystyle  |\lambda_0| \leq \max_{1 \leq j \leq n} \bigg\{1+|a_j|, 
\sum_{j=1}^{n} \sum_{k=1}^{m_j}||B_k^{(j)}||+ \sum_{i=0}^{m-1}||C_i|| \bigg\}$.

\item 
$ |\lambda_0| \leq 
\displaystyle \max_{\substack{\text{$1 \leq j \leq n$}\\ \text {$1 \leq i \leq m-1$}}} 
\bigg\{|a_j|+||B^{(j)}||, 1+||C_i||, ||C_0||+ n \bigg\}$, if each pole is of order $1$.\\


\ \ $\leq \displaystyle \max_{\substack{\text{$1 \leq j \leq n$}\\ 
\text{$1 \leq k \leq m_j$} \\ \text {$1 \leq i \leq m-1$}}} 
\bigg\{|a_j|+||B_k^{(j)}||, 1+|a_j|, 1+||C_i||, ||C_0||+ 
\sum_{j=1}^{n}m_j \bigg\}$, otherwise.

%\item $\displaystyle |\lambda_0| \leq \max_{\substack{1 \leq j \leq n \\ 1 \leq k \leq m_j \\
%1 \leq i \leq m-1}} \bigg\{|a_j|+||B_k^{(j)}||, 1+|a_j|, 1+||C_i||, ||C_0||+ 
%\sum_{j=1}^{n}m_j \bigg\}$.
\
\item $|\lambda_0| \leq \displaystyle \frac{\alpha+ \beta +\sqrt{(\alpha-\beta)^2 + 
(\gamma+\delta)^2}}{2}$,

where $\alpha = \displaystyle \max_{1 \leq j \leq n} \Big\{|a_j| + 
\cos\Big(\frac{\pi}{m_j+1}\Big) \Big\}$, $\beta = w(\mathcal{B}_0)$, $\gamma =
\sqrt{\displaystyle \sum_{j=1}^{n} m_j}$ and $\delta = \sqrt{ \displaystyle 
\sum_{j=1}^{n}\sum_{k=1}^{m_j} {||B_{k}^{(j)}||}^2}$ with $\mathcal{B}_0 = 
\begin{bmatrix}
0 & 1  & \cdots & 0\\
0 & 0 & \cdots & 0\\
\vdots  & \vdots & \ddots & \vdots\\
0 & 0 & \cdots & 1\\
||C_0|| & ||C_1|| & \cdots & ||C_{m-1}||\\ 
\end{bmatrix}$.
\end{enumerate}
\end{theorem}

\begin{proof}
Since $\lambda_0$ is an eigenvalue of $R(\lambda)$, by Theorem 
\ref{Thm-eigenvalue bounds-matrix rational function} $|\lambda_0| \leq R$, where 
$R$ is real root of the associated rational function $q(x)$ such that $|a_j| < R$ 
for $1 \leq j \leq n$.
Now $(1)$ and $(2)$ follow from Theorem \ref{Thm-bounds for zeros using norms}, 
and $(3)$ follow from Theorem \ref{Thm-bounds on zeros using numerical range-2}.
\end{proof}

\subsection{Special case} 

We now consider a rational function in which $p(\lambda)$ is a linear polynomial. 
That is, consider a rational function of the following form
\begin{equation}\label{Eqn- rational function}
r(\lambda) = \lambda - b_0 - \displaystyle \sum_{j=1}^{n} \sum_{k=1}^{m_j} 
\frac{b_k^{(j)}}{(\lambda-a_j)^k}.
\end{equation}

We then have the following bound on the zeros of $r(\lambda)$.

%\noindent
%Corresponding to $r(\lambda)$ the associated block square matrix $\mathcal{C}_r$ of 
%size $ 1+ \displaystyle \sum_{j=1}^{n} \frac{m_j(m_j+1)}{2}$ is  
%\begin{center}
%$\mathcal{C}_r= \begin{bmatrix}
%\mathcal{A}_n & 0 & \cdots & 0 & -\mathcal{F}_n \\
%0 & \mathcal{A}_{n-1} & \cdots & 0 & -\mathcal{F}_{n-1} \\
%\vdots & \vdots & \ddots & \vdots & \vdots \\
%0 & 0 & \cdots & \mathcal{A}_1 & -\mathcal{F}_1 \\
%\mathcal{B}_n & \mathcal{B}_{n-1} & \cdots & \mathcal{B}_1 & \mathcal{B}_0
%\end{bmatrix} $,
%\end{center}
%where for each $1 \leq j \leq n$, the blocks are given by: 
%
%\noindent 
%$\mathcal{A}_j = \text{diag}(A_{m_j}^{(j)},A_{m_j-1}^{(j)},\cdots,A_1^{(j)})$ is a 
%block diagonal matrix of size $\displaystyle \frac{m_j(m_j+1)}{2}$, with $A_k^{(j)} 
%=\begin{bmatrix}
%a_j & 1 & 0 & \cdots & 0 & 0 \\
%0 & a_j & 1 & \cdots & 0 & 0 \\
%\vdots & \vdots & \vdots & \ddots & \vdots & \vdots \\
%0 & 0 & 0 & \cdots & a_j & 1 \\
%0 & 0 & 0 & \cdots & 0 & a_j
%\end{bmatrix}_{k \times k}$, $1\leq k \leq m_j$.
%
%\
%
%\noindent
%$\mathcal{B}_j = \begin{bmatrix}
%B_{m_j}^{(j)} & B_{m_j-1}^{(j)} & \cdots & B_1^{(j)}
%\end{bmatrix}$ is a block row matrix of size $\displaystyle 1 \times 
%\frac{m_j(m_j+1)}{2}$, with $B_k^{(j)} = \begin{bmatrix}
%-b_{k}^{(j)} & 0 & \cdots & 0
%\end{bmatrix}_{1\times k}, 1 \leq k \leq m_j $, $\mathcal{B}_0 = \begin{bmatrix}
%b_0
%\end{bmatrix}$.
%
%\
%
%\noindent
%$\mathcal{F}_j = \begin{bmatrix}
%F_{m_j} \\
%F_{m_j-1} \\
%\vdots \\
%F_1
%\end{bmatrix}$ is block column matrix of size $\displaystyle \frac{m_j(m_j+1)}{1} 
%\times 1$ with $ F_k =\begin{bmatrix}
%0 \\
%\vdots \\
%0 \\
%1
%\end{bmatrix}_{k\times 1}, 1 \leq k \leq m_j$ for $j = 1,\cdots, n$.
%

\begin{theorem}\label{Thm-bounds for zeros}
Let $r(\lambda)$ be a rational function as given in Equation \eqref{Eqn- rational function}. 
If $\lambda_0$ is a zero of $r(\lambda)$, then $|\lambda_0| \leq 
\displaystyle \max_{1\leq j \leq n} \Big\{|a_j| + 
\cos\bigg(\frac{\pi}{m_j+1}\bigg), |b_0| \Big\} + \frac{1}{2} \Bigg( 
\sqrt{\sum _{j=1}^{n}m_j} + \sqrt{\sum_{j=1}^{n} \sum_{k=1}^{m_j}|b_k^{(j)}|^{2}} \Bigg)$.
\end{theorem}

\begin{proof}
By Corollary \ref{Cor-zeros of rational function}, $\lambda_0$ is an eigenvalue of 
$\mathcal{C}_r$. Therefore 
$|\lambda_0| \leq w(\mathcal{C}_r)$. We can split $\mathcal{C}_r$ as $\mathcal{C}_r 
= \mathcal{C}_1 + \mathcal{C}_2$, where $\mathcal{C}_1 = 
\text{diag}(\mathcal{A}_n,\mathcal{A}_{n-1}, \cdots , \mathcal{A}_1, \mathcal{B}_0)$  
and $\mathcal{C}_2 = \begin{bmatrix}
0 & 0 & \cdots & 0 & -\mathcal{F}_n \\
0 & 0 & \cdots & 0 & -\mathcal{F}_{n-1} \\
\vdots & \vdots & \ddots & \vdots & \vdots \\
0 & 0 & \cdots & 0 & -\mathcal{F}_1 \\
\mathcal{B}_n & \mathcal{B}_{n-1} & \cdots & \mathcal{B}_1 & 0
\end{bmatrix}$. Then we have $w(\mathcal{C}_r) \leq w(\mathcal{C}_1) + 
w(\mathcal{C}_2)$. It is easy to see that $w(\mathcal{C}_1) \leq \displaystyle 
\max_{1 \leq j \leq n} \Big\{|a_j| + \cos\Big(\frac{\pi}{m_j+1}\Big), |b_0| \Big\}$. 
Now consider $\mathcal{C}_2 = \begin{bmatrix}
0 & 0 & \cdots & 0 & -\mathcal{F}_n \\
0 & 0 & \cdots & 0 & -\mathcal{F}_{n-1} \\
\vdots & \vdots & \ddots & \vdots & \vdots \\
0 & 0 & \cdots & 0 & -\mathcal{F}_1 \\
\mathcal{B}_n & \mathcal{B}_{n-1} & \cdots & \mathcal{B}_1 & 0
\end{bmatrix} = \begin{bmatrix}
0 & \mathcal{F} \\
\mathcal{B} & 0
\end{bmatrix}$, where $\mathcal{F} = -\begin{bmatrix}
\mathcal{F}_n \\
\mathcal{F}_{n-1} \\
\vdots \\
\mathcal{F}_1
\end{bmatrix}$ and $ \mathcal{B} = \begin{bmatrix}
\mathcal{B}_n & \mathcal{B}_{n-1} & \cdots & \mathcal{B}_1
\end{bmatrix}$. By Lemma \ref{Lem-numerical radius-2} we have, $w(\mathcal{C}_2) \leq \displaystyle 
\frac{1}{2} (||\mathcal{F}||_2+ ||\mathcal{B}||_2) = \frac{1}{2} 
\Bigg( \sqrt{\sum _{j=1}^{n}m_j} + \sqrt{\sum_{j=1}^{n} 
\sum_{k=1}^{m_j}|b_k^{(j)}|^2} \Bigg)$. Therefore $|\lambda_0| \leq \displaystyle 
\max_{1 \leq j \leq n} \Big\{|a_j| + \cos\bigg(\frac{\pi}{m_j+1}\bigg), |b_0| \Big\}\\
 + \frac{1}{2} \Bigg( \sqrt{\sum _{j=1}^{n}m_j} + \sqrt{\sum_{j=1}^{n} 
\sum_{k=1}^{m_j}|b_k^{(j)}|^2} \Bigg)$. 
\end{proof}

\medskip
As a corollary we have the following result for matrix rational functions whose matrix 
polynomial part is linear.

\begin{corollary}\label{Cor-bound for eigenvalue-linear}
Let $R(\lambda) = I \lambda  - B_0 + \displaystyle \sum_{j=1}^{n} \sum_{k=1}^{m_j} 
\frac{B_k^{(j)}}{(\lambda-a_j)^k}$ be a $p\times p$ matrix rational function. If 
$\lambda_0$ is an eigenvalue of $R(\lambda)$, then $|\lambda_0| \leq \displaystyle  
\max_{1\leq j \leq n} \Big\{|a_j| + \cos\bigg(\frac{\pi}{m_j+1}\bigg), ||B_0|| \Big\} + 
\frac{1}{2} \Bigg( \sqrt{\sum _{j=1}^{n}m_j} + \sqrt{\sum_{j=1}^{n} 
\sum_{k=1}^{m_j}||B_k^{(j)}||^{2}} \Bigg)$.
\end{corollary}

\medskip
\section{Numerical results}\label{sec-4}\hspace*{\fill}
 
In this section, we determine the bounds on the eigenvalues of a REP 
and bounds on the zeros of a scalar rational function using the above 
proposed methods. Among several bounds given above for the eigenvalues of a general 
matrix rational function, it is obvious that the bound obtained using 
Theorem \ref{Thm-eigenvalue bounds-matrix rational function} is better. In general, 
we cannot compare other bounds, unless a particular example is considered. Similarly 
we cannot determine which method gives a better bound for zeros of scalar 
rational functions.      

\begin{example}\label{example-4.1}
Consider the following $3 \times 3$ matrix rational function, 
\begin{align*}
R(\lambda) = I \lambda^3+C_1 \lambda +C_0 +\displaystyle 
\frac{B_1}{(\lambda-1)^2} + \frac{B_2}{\lambda-2} 
\end{align*}
where $C_0 =\begin{bmatrix}
4 & 1 & 1 \\
1 & 4 & 1 \\
1 & 1 & 4
\end{bmatrix}$, $C_1 =\begin{bmatrix}
2 & 1 & 1 \\
1 & 2 & 1 \\
1 & 1 & 2
\end{bmatrix}$ and $B_1 = B_2 = I$. The polynomial part $P(\lambda)= 
I \lambda^3+C_1 \lambda +C_0$ is taken from \cite{Bani-Fuad-Rawan}. 
The maximum of the absolute values of eigenvalues of $R(\lambda)$ is $2.29$. 
The bounds obtained using Theorems \ref{Thm-eigenvalue bounds-matrix rational function} 
and \ref{Thm summary} are given in Table $1$ below.

\begin{table}[h!]
\centering
\begin{tabular}{l l l l l l}
\hline
Result & \hspace{5cm} & Bound \\ [0.5ex] 
\hline

Theorem 3.6 & \hspace{5cm} & 2.64\\
Theorem 3.7 (1) & \hspace{5cm} & 12.00\\
Theorem 3.7 (2) & \hspace{5cm} & 9.00\\
Theorem 3.7 (3) & \hspace{5cm} & 4.99\\[1ex] 
\hline 
\end{tabular}
\vspace{0.2cm}
\caption{Bounds obtained by Theorems $3.6$ and $3.7$}
\label{table:2}
\end{table} 
\end{example}

\begin{example}\label{example-4.2}
The following is a scalar rational function whose polynomial part 
is taken from \cite{Abdelkader-Faud-Soumia}
\begin{align*}
r(\lambda) = \lambda^5 + \lambda +4- \displaystyle \frac{1}{\lambda -1}+ 
\frac{2}{(\lambda -1)^2}+\frac{3}{\lambda -3}+\frac{4}{(\lambda -3)^2}-
\frac{1}{(\lambda -3)^2}.
\end{align*}
The bounds on the zeros of $r(\lambda)$ obtained using Theorems 
\ref{Thm-bounds for zeros using norms} and 
\ref{Thm-bounds on zeros using numerical range-2} are given in Table $2$ below. 
The maximum of the absolute values of zeros of $r(\lambda)$ is $3.12$.

\begin{table}[h!]
\centering
\begin{tabular}{l l l l l l}
\hline
Result & \hspace{5cm} & Bound \\ [0.5ex] 
\hline

Theorem 3.3 (1) & \hspace{5cm} & 16.00\\
Theorem 3.3 (2) & \hspace{5cm} & 9.00\\
Theorem 3.4 & \hspace{5cm} & 6.96\\[1ex] 
\hline
\end{tabular}
\vspace{0.2cm}
\caption{Bounds obtained by Theorems \ref{Thm-bounds for zeros using norms} and 
\ref{Thm-bounds on zeros using numerical range-2}}
\label{table:2}
\end{table} 

Note that the roots of $r(\lambda)$ can be found by converting it to a 
rational function of the form $\displaystyle \frac{p(\lambda)}{q(\lambda)}$, 
where $p(\lambda)$ and $q(\lambda)$ are scalar polynomials which are coprime. 
Then the roots of $r(\lambda)$ and $p(\lambda)$ are the same. Now one can derive 
bounds on roots of $p(\lambda)$ using various techniques 
(see \cite{Abdelkader-Faud-Soumia}). In practice, however, it is difficult 
to determine the coefficients of $p(\lambda)$ and coefficients might be very 
large. For the above example, the bound obtained using Theorem 
\ref{Thm- eigenvalues- matrix rational function} of 
\cite{Abdelkader-Faud-Soumia} is $277.47$, which is very large when compared 
to the bounds obtained in Table $2$. 

\end{example}

\medskip
\noindent
{\bf Declarations:} The authors declare that there is no conflict of interest in this work.

\noindent
{\bf Acknowledgements:} Pallavi .B and Shrinath Hadimani are grateful for the financial 
assistance provided by the Council of Scientific and Industrial Research (CSIR) and the 
University Grants Commission (UGC), both departments of the Indian government.

\bibliographystyle{amsplain}
\begin{thebibliography}{99}

\bibitem{Abu-Kittaneh}
A. Abu-Omar and F. Kittaneh, \emph{Estimates for the numerical radius and the spectral 
radius of the {F}robenius companion matrix and bounds for the zeros of polynomials}, 
Annals of Functional Analysis, 5(1) (2014), 56-62.

\bibitem{Amparan-Dopico}
A. Amparan, F. M. Dopico, S. Marcaida, and I. Zaballa, \emph{On minimal bases and indices 
of rational matrices and their linearizations}, Linear Algebra and its Applications, 
623 (2021), 14-67.

\bibitem{R Bhatia}
R. Bhatia, \emph{Matrix Analysis}, Springer-Verlag, New York, 1997.

\bibitem{Bani-Fuad-Rawan}
W. Bani-Domi, F. Kittaneh and R. Mustafa, \emph{Bounds for the eigenvalues of 
matrix polynomials with commuting coefficients},  Results in Mathematics, 
78(3) (2023), 1422-6383

\bibitem{Dopico}
F. M. Dopico, S. Marcaida, M. C. Quintana and P. Van Dooren, \emph{Local linearizations 
of rational matrices with application to rational approximations of nonlinear eigenvalue 
problems}, Linear Algebra and its Applications, 604 (2020), 441-475.

\bibitem{Abdelkader-Faud-Soumia}
A. Frakis, F. Kittaneh and S. Soltani, \emph{New numerical radius inequalities 
for operator matrices and a bound for the zeros of polynomials}, Advances in 
Operator Theory, 8(1) (2023), 2662-2009.

\bibitem{Gustafson-Rao}
K. E. Gustafson and D. K. M. Rao, \emph{Numerical range-The field of values of linear 
operators and matrices}, Springer-Verlag, New York, (1997).

%\bibitem{Hirzallah-Kittaneh}
%O. Hirzallah, F. Kittaneh and K. Shebrawi, \emph{Numerical radius inequalities for 
%certain {$2\times2$} operator matrices}, Integral Equations and Operator Theory, 71(1) 
%(2011), 129-147.

\bibitem{Kailath}
T. Kailath, \emph{Linear systems}, Prentice-Hall Information and System Sciences Series, 1980.

\bibitem{Mehrmann-Voss}
V. Mehrmann and H. Voss, \emph{Nonlinear eigenvalue problems: a challenge for modern 
eigenvalue methods}, Mitteilungen der Gesellschaft f\"{u}r Angewandte Mathematik und 
Mechanik, 27(2) (2004), 121-152

\bibitem{Pallavi}
Pallavi. B, S. Hadimani and S. Jayaraman, \emph{Bounds on the eigenvalues of matrix 
rational functions}, arXiv:2302.02894.

\bibitem{Ponnusamy}
S. Ponnusamy, \emph{Foundations of complex analysis}, Narosa Publishing House, New Delhi, 
(1995).

\bibitem{Rosenbrock}
H. H. Rosenbrock, \emph{State-space and Multivariable Theory}, John Wiley and Sons, 
New York, 1970.
 	  	
\bibitem{Saad-Guide-Meidlar}
Y. Saad, M. El-Guide and A. Miedlar, \emph{A rational approximation method for nonlinear 
eigenvalue problem}, arXiv:1901.01188.

\bibitem{Su-Bai}
Y. Su and Z. Bai, \emph{Solving rational eigenvalue problems via linearization}, SIAM Journal 
on Matrix Analysis and Applications, 32(1) (2011), 201-216.	
\end{thebibliography}

%$|\lambda_0|^{-(m-1)} ||C_1||$
\end{document}